\section{Proofs and certificates for the foundations}
\label{fd:appendix}

\subsection{Cone containment, transversality, and boundary failures}

\begin{proof}[Proof of \cref{fd:cone-attainment}]
\label{fd:cone-proof}
If $S$ is empty, the whole space gives an infeasible intersection cut,
which dominates any valid inequality. Otherwise work in
$V=\operatorname{span}P$. For a valid $c\ge0$, the dual polyhedron
$D_c=\{y\in V:P^Ty\le c\}$ is nonempty and pointed: its lineality is
$\ker(P^T)\cap V=\{0\}$. It has finitely many vertices
$y^1,\ldots,y^m$. The fibre LP has an attained finite minimum for every
$v\in\cone(P)$, and duality and the existence of a vertex optimum give
\[
 \phi(v)=\min_{\lambda\ge0:P\lambda=v}c^T\lambda
        =\max_{1\le i\le m}(y^i)^Tv.
\]
Define the full-dimensional polyhedron
\[
 C=\{\sbar+v+u:v\in V,\ u\in V^\perp,\ (y^i)^Tv\le1
       \text{ for every }i\}.
\]
It contains a neighbourhood of the apex. A point of $S$ has $u=0$ and
$v\in\cone(P)$, and validity gives $\phi(v)\ge1$; hence it is not in
$\intt C$. Along ray $p_j$, every $0\le t\le1/c_j$ lies in $C$, and
the entire ray lies there if $c_j=0$. Thus $a_j(C)\le c_j$, proving
domination.

For $0<\zK(\rc)<\infty$, apply this result to $c=\rc/\zK(\rc)$.
If $\zK=0$, every free neighbourhood has value zero by validity.
If $\zK=+\infty$, cone containment forces $S=\varnothing$, so the
whole-space set attains that value. This proves the stated attainment
claim even when $\rc$ has zero entries.

When $\cone(P)=\R^k$ and $0<z=\zK(\rc)<\infty$, the cost image
$V_z=\sbar+P\{\lambda\ge0:\rc^T\lambda\le z\}$ itself is a closed
polyhedron, a free neighbourhood, and an attaining set.
To see interior membership, represent each $\pm e_i$ by nonnegative
ray vectors; their maximum finite cost bounds the cost of any
sufficiently small vector. If $s\in\intt V_z$, extend the segment
from $\sbar$ through $s$ slightly within $V_z$ and rescale its
ray representation to obtain a representation of $s$ with cost
strictly below $z$. Thus $s$ cannot be feasible. Its original
ray steps are at least $z/\omega_j$, or infinite for zero costs.
\end{proof}

\begin{proof}[Proof of \cref{fd:attainment}]
\label{fd:attainment-proof}
Put $T=T_z$ and $M=T\cap S$. Every representation of a point of $M$
with objective at most $z$ has objective exactly $z$. Moreover $q=0$
on $M$: a point with $q<0$ could be scaled slightly toward the apex
to produce a feasible ray vector of smaller cost. In particular,
$q\ge0$ on $T$. The contact set $M$ is compact. Continuity and the
strict hypothesis give a neighbourhood $U$ of $M$, a neighbourhood
$W$ of $\sbar$, and $\eta>0$ such that
\[
 \nabla q(\xi)^T(b-t)\ge\eta
\]
whenever $t,\xi\in U$ are sufficiently close to each other and $b\in W$.
For example, first take a positive minimum of
$\nabla q(t)^T(\sbar-t)$ on $M$, and then use uniform continuity on
compact neighbourhoods.

Suppose for arbitrarily small $\delta_n\downarrow0$ there are
$p_n\in C_{\delta_n}\cap S$ outside $M$. Write
$p_n=\theta_n b_n+(1-\theta_n)t_n$, with $b_n$ in the apex ball,
$t_n\in T$, and $0\le\theta_n\le1$. A subsequence has
$p_n\to p\in M$, $t_n\to t\in T$, and $\theta_n\to\theta$.
If $\theta>0$, a ray representation of $t$ with cost at most $z$
gives a representation of
$p=\theta\sbar+(1-\theta)t$ with cost at most $(1-\theta)z<z$.
This contradicts $p\in S$ without requiring $P$ to be injective.
Thus $\theta=0$ and $t_n\to p\in M$.
When $\theta_n=0$, $p_n=t_n$ already belongs to $M$, a contradiction.
Otherwise the segment mean value theorem gives, for large $n$,
\[
 q(p_n)=q(t_n)+\theta_n\nabla q(\xi_n)^T(b_n-t_n)
        \ge\theta_n\eta>0,
\]
again a contradiction. Therefore $C_\delta\cap S=M$ for all
sufficiently small $\delta$.
If a point of $M$ were interior to $C_\delta$, it would be an interior
local minimum of $q$, so its gradient would vanish, contrary to the
strict hypothesis. Thus $C_\delta$ is free. It contains $T_z$, and
validity now gives $z_{C_\delta}=z$.
\end{proof}

\begin{proof}[Proof of \cref{fd:boundary-examples}]
\label{fd:boundary-proof}
For $q=(x+1)y+1$, the positive orthant has $q\ge1$, so the corner is
infeasible. Any free neighbourhood contains $B(0,\delta)$. If its
first ray step were infinite, $e_1$ would be a recession direction,
and $(2/\delta,-\delta/2)$ would be interior to the set. But its
quadratic value is $-\delta/2<0$. The first step is therefore finite,
so the zero-cost first ray makes every cut value zero.
For $q=(x+1)y+1-y^2$, the same point has value
$-\delta/2-\delta^2/4<0$. Orthant feasibility requires
$y^2-y-1\ge0$, so $y\ge(1+\sqrt5)/2$, attained at $x=0$.

For $q=x-x^2+(1-y)^2$, any orthant point of cost at most one has
$0\le x\le1$ and $q\ge0$. Equality forces $(x,y)=(0,1)$;
this point is feasible, so the corner value is one with a unique
minimizer. If a cut had value at least one, its second step would be
at least one and $(0,1)$ would lie in its defining free set.
For sufficiently small $s>0$, $(-s,0)$ lies in the apex ball.
The point
\[
 (1-s)(0,1)+s(-s,0)=(-s^2,1-s)
\]
is then interior to the convex set, while its quadratic value is
$-s^4<0$. This proves non-attainment.
\end{proof}

\subsection{Algorithmic details}

\begin{proof}[Proof of \cref{fd:algorithms}]
\label{fd:algorithm-proof}
By \cref{fd:rank-support,fd:inertia}, a feasible optimum uses at most
$r$ rays. Thus $\zK\le t$ exactly when, for one of the $O(N^r)$ supports
$J$, the fixed-variable existential formula
\[
 \exists\mu\in\R^J:\quad
 \mu\ge0,\quad \rc_J^T\mu\le t,\quad q(\sbar+P_J\mu)\le0
\]
holds. Its number of variables and polynomial degrees are bounded
for fixed $r$. Fixed-variable real-algebraic decision therefore has
polynomial bit complexity \citep[Theorem~1.1]{Renegar1992PartIII}.
The same formulas without the objective
bound detect infeasibility.

For completeness, a finite positive value has a polynomial-bit initial
bracket. On each fixed-size support, real quantifier elimination with
the objective threshold left free yields univariate polynomials of
bounded degree and coefficient bit size polynomial in the input
size. A finite endpoint is a root of one of those polynomials.
Standard root-size bounds place a positive endpoint between
$2^{-B}$ and $2^B$ for a polynomially bounded $B$. There are only
polynomially many supports. Bisection on a common such bracket,
using the decision test, yields relative error $\epsilon$ with
polynomially many calls in the input size and $\log(1/\epsilon)$.

For the normal-fan refinement, assume $\rank P=k$ and consider
$D_\rc=\{y:P^Ty\le\rc\}$. It is pointed and contains zero in its
interior. Fibre duality gives
\[
 \phi(v)=\min_{\lambda\ge0:P\lambda=v}\rc^T\lambda
        =\max_{y\in D_\rc}y^Tv
 \quad(v\in\cone(P)).
\]
For a vertex $y$, $\phi(v)=y^Tv$ on its normal cone, generated by
the active columns $p_j$.
Perturb the right-hand side by a sufficiently small generic
lexicographic perturbation, keeping every entry positive, to obtain
a simple polyhedron. Each active independent $k$-set $J$ gives a
vertex depending continuously on that perturbation. There are
finitely many such sets. If its limit were infeasible for the original
polyhedron, it would remain infeasible under small perturbations.
Thus every perturbed vertex limits to an original vertex, and its
normal cone $\cone(P_J)$ lies in that original normal cone.
The perturbed vertex cones cover $\cone(P)$ because every bounded
dual objective there has a vertex optimum.
On such a cell,
$\phi(P_J\mu)=y^TP_J\mu=\rc_J^T\mu$, with the original costs.
The cell problems therefore preserve the original value; their
optima use at most $r$ generators by \cref{fd:inertia}.

A pointed unbounded polyhedron can be truncated by one additional
facet while keeping all its vertices. The upper bound theorem for
the resulting polytope gives $O(N^{\lfloor k/2\rfloor})$ cells
\citep{McMullen1970}.
For $k=3$, a bounded simple polyhedron with at most $N$ facets has at
most $2N-4$ vertices. In the unbounded case the truncation has at most
$N+1$ facets and $2N-2$ vertices, at least three on the new facet,
so the same bound for original vertices is valid. The normal-fan
edge graph embeds on the sphere, has at most $N$ ray-direction
vertices, and is simple after zero and duplicate directions are
removed. Its planar bound is $3N-6$ edges.
\end{proof}

\begin{remark}[Generic stationary points]\label{fd:generic-kkt}
On a face $F$ of a support subproblem, write
$g(\mu)=\mu^TG_F\mu+2m_F^T\mu+g_0$.
For a regular stationary point with nonsingular $G_F$,
\[
 \mu=-G_F^{-1}(m_F+\tau\rc_F),\qquad
 \tau^2=\frac{m_F^TG_F^{-1}m_F-g_0}
                  {\rc_F^TG_F^{-1}\rc_F},\qquad \tau>0.
\]
The denominator must be nonzero, and candidates must satisfy the
coordinate and active-constraint conditions. Enumerating faces gives
the generic closed-form solution. Singular faces, abnormal stationary
points, and zero denominators require the exact algebraic decision
method or, with at most two rays, the next formula.
\end{remark}

\begin{proof}[Proof of \cref{fd:two-ray}]
\label{fd:two-ray-proof}
Every nonzero nonnegative two-ray vector is $\tau\nu(\theta)$, with
cost $\tau>0$. Substituting $u=1/\tau$ gives
$g_0u^2+b(\theta)u+a(\theta)\le0$. Since $g_0>0$, the greatest feasible
positive reciprocal radius is its larger root $u_+(\theta)$.
The compact discriminant-feasible set gives the stated maximum;
a positive larger root itself is feasible.

Away from $\Delta=0$,
$u_+'=(-b'+\Delta'/(2\sqrt\Delta))/(2g_0)$, so interior stationary points
satisfy $\Delta'=2b'\sqrt\Delta$ and hence
$\Delta'^2-4b'^2\Delta=0$. This polynomial has degree at most two, because
$b$ is affine and $\Delta$ is quadratic. Interval boundaries occur at
zero or one or at a root of $\Delta$. Squaring may add candidates, but
evaluating the actual larger root discards no valid maximum.
Including roots of $\Delta'$ also handles the $b'=0$ case directly.
If the stationary polynomial vanishes identically and $b'\ne0$,
$\Delta$ is a square of an affine function, so $\sqrt\Delta$ and $u_+$ are
piecewise affine; maxima occur at boundaries or kinks $\Delta=0$.
If $b'=0$, identically vanishing stationarity means $\Delta$ is constant,
and the root is constant. If $\Delta\equiv0$, the root is affine.
All exceptional cases are therefore covered by finitely many
endpoint and quadratic-root evaluations. If no positive larger root
exists, the two-ray subproblem is infeasible.
\end{proof}

\begin{proof}[Bilinear attainment qualification]
\label{fd:bilinear-attainment-proof}
The gradient of $q=w-xy$ never vanishes. Suppose a minimizing
contact $t$ fails \cref{fd:transverse-condition}. Choose an optimal
representation of that same point with independent supported columns
using its fibre LP. If the supported gradient were nonzero, the
regular KKT system on its positive support would give a positive
objective multiplier and thus positive transversality.
Therefore the supported gradient vanishes. The abnormal case of
the inertia proof bounds its support by $\rho-1=1$.
The contact is consequently a single ray's first feasible point,
and its directional gradient is zero. If no such minimizing contact
exists, all minimizing contacts satisfy the transversality hypothesis,
so \cref{fd:attainment} applies.
\end{proof}

\begin{proof}[Proof of \cref{fd:hardness}]
\label{fd:hardness-proof}
For a graph $G$ on $k$ vertices, let $A_G$ have zero diagonal and
entries one on its edges, and let $\omega(G)$ denote its clique number. Set
$q(s)=\mu-s^TA_Gs$, $P=I$, $\sbar=0$, and $\rc=\mathbf1$.
Writing $s=\tau\nu$, $\nu\ge0$, $\mathbf1^T\nu=1$, the Motzkin--Straus
identity \citep{MotzkinStraus1965} gives
\[
 \max_\nu\nu^TA_G\nu=1-\frac1{\omega(G)},\qquad
 \zK=\sqrt{\frac{\mu\,\omega(G)}{\omega(G)-1}}.
\]
An edgeless graph has infinite value and can be recognized directly.
For a rational integer threshold $\omega_0\ge2$, use
$\mu=\omega_0(\omega_0-1)$. Then
$\zK\le\omega_0$ exactly when $\omega(G)\ge\omega_0$, proving
rational-threshold NP-hardness.
With $\mu=1$ put $f(h)=\sqrt{h/(h-1)}$. For consecutive clique sizes,
\[
 \frac{f(h)}{f(h+1)}
 =\left(1-\frac1{h^2}\right)^{-1/2}
 \ge1+\frac1{2k^2}.
\]
For $\delta=1/(5k^2)$ and $k\ge2$,
$(1+\delta)/(1-\delta)<1+1/(2k^2)$.
The relative-error intervals around all possible finite values
are disjoint, so a $\delta$-approximation determines the clique
number. The matrix $-A_G$ is indefinite whenever $G$ has an edge.
\end{proof}

\begin{remark}[A specified efficacy optimization]\label{fd:efficacy}
In ray-coordinate Euclidean norm, a normalized valid cut
$a^T\lambda\ge1$ has efficacy $1/\norm a$. Thus maximize efficacy by
minimizing $\norm a$ over
$\mathcal V=\{a\ge0:a^T\lambda\ge1\text{ on }X\}$.
For fixed $r$, the support-value identity extends from positive to
nonnegative costs by perturbation: for any $a\ge0$,
$\zK(a+\epsilon\mathbf1)\downarrow\zK(a)$, and the finite union of
supports of size at most $r$ has the same infimum. The fixed-variable
algebraic procedure can therefore test $\exists\lambda\in X:
a^T\lambda<1$ and return a violating point. Such a point gives the
separating inequality $\lambda^Tc\ge1$ for all $c\in\mathcal V$.
Algebraic sample points can be approximated to the accuracy of a
weak separation oracle.

If $X$ is nonempty, write $z_1=\zK(\mathbf1)>0$.
The vector $(2/z_1)\mathbf1$ lies in $\mathcal V$ with a ball of
radius $1/z_1$: for every $\lambda\in X$,
$\norm\lambda\le\mathbf1^T\lambda$ and
$(2/z_1)\mathbf1^T\lambda-1\ge\norm\lambda/z_1$.
A feasible vector of norm $2\sqrt N/z_1$ bounds the optimum.
Set $a_0=(2/z_1)\mathbf1$ and $R=4\sqrt N/z_1$, and introduce
the bounded epigraph
\[
 \mathcal E=\{(a,t):a\in\mathcal V,\ \norm a\le t\le2R\}.
\]
Its linear objective $t$ has the same minimum as the norm problem.
It contains the ball of radius $1/(2z_1)$ about $(a_0,R)$:
the $a$ component stays inside the established feasible ball,
the norm constraint has initial slack $R/2$, and a joint displacement
of that radius changes $\norm a-t$ by at most $\sqrt2/(2z_1)$.
The upper constraint has slack $R$. The whole epigraph lies in
the ball of radius $2\sqrt2R$ about zero.
The weak separation oracle for $\mathcal V$, together with
elementary separation for the norm and linear constraints,
therefore supplies a weak separation oracle for this bounded body.
Weak separation-to-linear-optimization equivalence
\citep[Corollary~4.2.7]{GrotschelLovaszSchrijver1988}
minimizes $t$ to prescribed accuracy in polynomial time for fixed $r$.
For rational bit encoding, choose a rational $b$ with
$0<b\le z_1\le2b$, replace $a_0$ by $(2/b)\mathbf1$ and $R$ by
$4\lceil\sqrt N\rceil/b$, and use radius $1/(2b)$ for the inner
ball and $3R$ for the outer ball. The same estimates apply;
the preceding algebraic size bound controls their bit sizes.
This concerns a stated ray-coordinate metric. An
original-variable metric must be supplied through its linear map;
it is not implicit in this assertion. If $X$ is empty, the
normalized-cut efficacy problem degenerates and is excluded.
The efficacy supremum of genuine intersection cuts has this same
value: perturb a valid nonnegative vector to be positive, and
apply \cref{fd:single-cut} at target $1-\epsilon$ to obtain cut
coefficients bounded by the perturbed vector divided by
$1-\epsilon$. Let both perturbations vanish. Attainment by a
genuine intersection cut does not follow.
\end{remark}

\subsection{The two polytope closures}

\begin{proof}[Proof of \cref{fd:nonlocal-hull,fd:rank-one-example}]
\label{fd:polytope-proof}
Every set in the intersection defining the nonlocal closure contains
$\mathcal P\cap S$. If this intersection is nonempty, its convex hull
is compact. A point $p$ of $\mathcal P$ outside that hull can be
strictly separated by an inequality $a^Ts\ge b$. Choose
$a^Tp<b'<b$. The compact set
$\mathcal P\cap\{a^Ts\le b'\}$ misses closed $S$.
Its sufficiently small ball enlargement is free and contains that
compact set in its interior. All remaining points of $\mathcal P$
have $a^Ts>b'$, so their convex hull excludes $p$.
If $\mathcal P\cap S$ is empty, a small enlargement of all of
$\mathcal P$ is free and removes every point.

For the triangle example, the closed unit disk is the unique maximal
full-dimensional free set: every free interior lies in the open disk,
so every full-dimensional closed free set lies in its closure.
The top vertex belongs to $S$ and yields no cut. At $(-1/2,0)$ the
basis rays point toward $(0,2)$ and $(1/2,0)$. The first exit is
$p_1=(-1/2,0)+t((0,2)-(-1/2,0))$, where
$17t^2-2t-3=0$ gives the displayed positive root.
The second exit is $(1,0)$, beyond the triangle. The disk cut is
the line through these two exits and dominates every other cut
at that basis. The other lower vertex gives its reflection.
Their intersection is
\[
 (0,\,2t/(3/2-t/2))
 =\left(0,\frac{1+\sqrt{13}}6\right).
\]
The resulting polygon has vertices $(0,2),p_1,v,p_2$.
The feasible part of the triangle lies above the chord $p_1p_2$:
below height $2t$, its horizontal sections lie strictly inside the
disk. All three vertices of the triangle above that chord are
feasible, so the hull is precisely
$\conv\{(0,2),p_1,p_2\}$.
Since $(1+\sqrt{13})/6<2t$, the first closure is strictly larger.
At its new vertex $v$, the two active cuts give rays through
$p_1$ and $p_2$, so the next disk cut is their chord and removes
the remaining gap.
\end{proof}

\subsection{Lorentz classification and the point rule}

\begin{proof}[Proof of \cref{fd:lorentz}]
\label{fd:lorentz-proof}
First let $m=1$. The forbidden homogeneous set is the union of
the two closed convex Lorentz cones
$Q_+=\{\norm{x}\le y\}$ and $Q_-=\{\norm{x}\le-y\}$.
Separate the interior of a full-dimensional free convex set from
each cone. Because the cones contain zero and are homogeneous,
the resulting halfspaces can have zero right-hand side and take
the form
\[
 v^Tx\ge y,\qquad w^Tx\ge-y,\qquad \norm v,\norm w\le1.
\]
Indeed a nonzero functional nonpositive on $Q_+$ has the form
$(v,-1)$ after positive normalization, with $\norm v\le1$;
the other cone gives the second inequality.
Their intersection contains the given free set and is free.
Full dimension requires $v+w\ne0$.

With the Lorentz form $B(\xi,\eta)=\xi^T
\diag(I_n,-1)\eta$, put $A=(v,-1)$ and $B_0=(w,1)$.
The quadratic values of $A$ and $B_0$ are nonpositive, whereas
the value at $A+B_0$ is $\norm{v+w}^2>0$.
Consequently $A+tB_0$ and $B_0+sA$ are null for some
$0\le t,s<1$. Their halfspaces still contain the original
intersection, since the normals are nonnegative combinations.
Their last entries are respectively negative and positive;
normalizing those entries gives the unit-direction pair in
\cref{fd:lorentz-pair}. It is full dimensional, hence
$\gamma_1\ne-\gamma_2$.

Every such pair is maximal. Let
$k_1=(\gamma_1,1)$, $k_2=(\gamma_2,-1)$, and
$h=1+\gamma_1^T\gamma_2>0$. At $k_1$ the first inequality is tight
and the second has slack $h$. If an added point $p$ violates the
first inequality, then
$(k_1-\theta p)/(1-\theta)$ lies in the interior of the pair for
sufficiently small $\theta>0$. Hence $k_1$ becomes an interior
point of the enlargement and is forbidden. The second inequality
has the same argument. Thus a maximal original set equals the pair.

To construct its orbit from a fixed unit $\ell$, set
$d=\sqrt{2/h}$,
$e'=d(k_1+k_2)/2$, and $\tau'=d(k_1-k_2)/2$.
These are Lorentz-orthonormal, with $e'$ spacelike and $\tau'$
future timelike. Map $(\ell,0)$ to $e'$ and $(0,1)$ to $\tau'$,
then extend by an isometry on their positive definite orthogonal
complements. Their dimension is $n-1\ge1$, so choose the extension
with determinant one. The resulting map is in $SO^+(n,1)$ and
sends $(\ell,\pm1)$ to $dk_{1,2}$. Lorentz orthogonality identifies
the corresponding supporting hyperplanes and their orientations,
so it sends $C_\ell$ to the target pair.

For $n=1$, the complement of $Q_{1,m}$ has two connected components,
$x>\norm y$ and $-x>\norm y$. The connected interior of any
full-dimensional convex free set lies in one component; its
closure lies in the corresponding closed cone. These cones are
free and maximal, proving the classification.

For an affine slice $H=\{u:\ell(u)=-1\}$ of a homogeneous form,
let $K$ be a full-dimensional free convex set relative to $H$.
The interior of its positive cone is
$\{rk:r>0,\ k\in\relint_H K\}$: the map $(r,k)\mapsto rk$
has inverse $u\mapsto(-\ell(u),u/[-\ell(u)])$ on $\{\ell<0\}$.
Homogeneity makes its closed convex cone free. The classification
just proved supplies a maximal homogeneous free set containing it.
Because the relative interior of $K$ lies in the homogeneous
interior, slicing that set back is free and contains $K$;
maximality of $K$ gives equality.

Finally, radical coordinates require a lineality argument.
If $S'+V=S'$ for a linear subspace $V$, any full-dimensional maximal
free set $C$ has $C+V=C$, since
$\intt\cl(C+V)=\intt C+V$ misses $S'$. Write the diagonalized slice
as $a^Tx+d^Ty+h^Tz=-1$, with $z$ radical. If $h=0$, this gives
the product with all radical coordinates. If $h\ne0$, the affine
bijection to the nondegenerate coordinates and $h^\perp$ is
\[
 (x,y,z_\perp)\longmapsto
 \left(x,y,z_\perp-\frac{1+a^Tx+d^Ty}{\norm h^2}h\right).
\]
The feasible set becomes $Q_{n,m}\times h^\perp$. Lineality and
maximality identify every free set as a product with a maximal
homogeneous free set. Equivalently, include the entire radical
product before imposing the original slice.
\end{proof}

\begin{proof}[Proof of \cref{fd:point-rule}]
\label{fd:point-rule-proof}
In transformed coordinates, the apex is
\[
 (\norm{x'}\ell,y')
 =\frac{\norm{x'}+y'}2(\ell,1)
  +\frac{\norm{x'}-y'}2(\ell,-1).
\]
Both coefficients are positive. Transformation gives
\cref{fd:point-span}. Conversely, the explicit Lorentz map in
the preceding proof sends $(\ell,\pm1)$ to $d k_{1,2}$, so its
inverse sends an apex satisfying \cref{fd:point-span} to
$((a+b)\ell/d,(a-b)/d)$. The normalized spatial direction is
exactly $\ell$, proving the converse.
For $n=2$, the identity $\norm{x-a\gamma_1}=b=a-y$ yields
the displayed expression for $a$, and then $b,\gamma_2$.
For any unit $\gamma_1$ with $\gamma_1^Tx-y>0$, the formula
gives $a>0$ and
$2(a-y)(\gamma_1^Tx-y)=\norm{x-y\gamma_1}^2>0$,
so $b>0$. This establishes the claimed parameterization and
its domain. Multiplying the two tangency rays to make their
slice coordinates one turns a positive apex combination into
a convex combination precisely when both scaling signs are positive.
\end{proof}

\begin{example}[A maximal set with only asymptotic contacts]
\label{fd:asymptotic-example}
For $S=\{xy\ge1\}$, let $K=\{x\ge0,y\le0\}$.
If $p\notin K$ has $p_x<0$, choose a small $\epsilon>0$ and large
$T>0$ so that the midpoint of $p$ and
$(\epsilon,-T)\in\intt K$ has two negative coordinates and product
greater than one. This midpoint lies in the interior of
$\conv(K\cup\{p\})$. If $p_y>0$, use $(T,-\epsilon)$.
Thus $K$ is maximal free, although $K\cap S=\varnothing$.

The homogenized form is $t^2-xy$, with signature $(2,1)$.
The two null tangency rays of $K$ are $(1,0,0)$ and $(0,-1,0)$.
Their span has $t=0$, so no affine apex with $t=1$ satisfies the
point-rule span condition. More generally any plain point-rule
slice contains a finite feasible contact: its apex is a positive
combination of its two null tangency rays and has slice coordinate
one, so at least one ray has positive slice coordinate and can
be scaled to a finite contact. Every enlargement of that set
still contains this contact. Since $K$ has none, the exclusion
also holds for the affine enlargement procedure of
\citet{MPS2022}, even if the apex is allowed to vary.
\end{example}

\begin{proof}[Proof of \cref{fd:fixed-rule}]
\label{fd:fixed-rule-proof}
The two point-rule steps are $x_0+1/x_0$ and $\sqrt{1+x_0^2}$,
by substitution along the two rays. The strip's steps are
$+\infty$ and one, giving its value one.
It is maximal free: it is invariant under horizontal translation,
and adding a point with $|y|>1$ makes a horizontal interior band
meet $S$. Corner minimization reduces to
\[
 \min_{\lambda_1\ge0}
 \epsilon\lambda_1+\sqrt{(x_0-\lambda_1)^2+1}.
\]
The stationary point has
$x_0-\lambda_1=\epsilon/\sqrt{1-\epsilon^2}$ and is feasible
under the stated hypothesis. Convexity gives the displayed value.
The asymptotic ratios follow on substituting $\epsilon=x_0^{-2}$.
\end{proof}

\begin{proposition}[The constant-direction closure on pointed corners]
\label{fd:constant-family}
For $S=\{y^2\ge x^2+1\}$, there are pointed two-ray corners with
positive costs for which the closure of all constant-direction
sets has value tending to zero, while the corner value and one
Lorentz-transformed split's value tend to one.
\end{proposition}
\begin{proof}
For $a>1$, put $r=\sqrt{1+a^2}$, $h=\sqrt{1+2a^2}$,
$n=(r,-a)/h$, $\eta=a^{-3}$, and use apex zero, costs $(1,1)$,
and rays $r_1=(a,r)$, $r_2=-r_1+\eta n$.
Their determinant is $-\eta h\ne0$, so the corner is pointed.
The boost $(X,Y)=(rx-ay,ry-ax)$ preserves $y^2-x^2$ and sends
the rays to $(0,1)$ and $(\eta h,-1-\eta c)$, where $c=2ar/h$.
For ray cost $\tau=\lambda_1+\lambda_2$, $|Y|\le(1+\eta c)\tau$;
feasibility requires $|Y|\ge1$. Thus
\[
 (1+\eta c)^{-1}\le\zK\le1.
\]
The free oblique split $|Y|\le1$ has value $(1+\eta c)^{-1}$,
which tends to one.

All constant-direction sets containing zero are
$|y|\le(ux+1)/\sqrt{1+u^2}$, $u\in\R$.
Both rays are timelike for this choice of $\eta$:
the second has squared Lorentz length
$1+2\eta c-\eta^2/h^2>0$.
Their cut coefficients are positive and have sum
\[
 a_1(u)+a_2(u)
 =(2r+\eta a/h)\sqrt{1+u^2}-(\eta r/h)u.
\]
Its minimum is
$L=\sqrt{4r^2+4\eta ar/h-\eta^2/h^2}$.
Consequently $(1/L,1/L)$ satisfies every family cut, and the
closure value is at most $2/L\to0$. This is an analytic pointed
perturbation, not an inference from the exactly opposite-ray limit.
\end{proof}

\subsection{The determinant orbit and its completion}

\begin{proof}[Proof of \cref{fd:orbit}]
\label{fd:orbit-proof}
A two-dimensional linear space all of whose matrices have rank
at most one has either a common left factor or a common right
factor. Indeed two rank-one matrices with independent left and
right factors have a rank-two sum. These two families of planes
are the two rulings of the determinant-zero cone.
A determinant-preserving bijection preserves or interchanges the
rulings: planes within one ruling meet only at zero, whereas
planes from opposite rulings meet in a rank-one line.
If it preserves them, select images of the four matrix units.
Their left and right factors give independent bases $a_1,a_2$
and $b_1,b_2$, with scalar coefficients $c_{ij}$.
The image of $(e_1+e_2)(e_1+e_2)^T$ has rank one, so
$c_{11}c_{22}=c_{12}c_{21}$. These nonzero coefficients split into
row and column factors, giving $M\mapsto AMB^T$.
If the rulings are interchanged, compose with transposition.
Determinant preservation then imposes $\det A\det B=1$.

With $J=\bigl(\begin{smallmatrix}0&1\\-1&0\end{smallmatrix}\bigr)$,
$K=\diag(1,-1)$, and
$L=\bigl(\begin{smallmatrix}0&1\\1&0\end{smallmatrix}\bigr)$, the map
\[
 \Phi(x_1,x_2,y_1,y_2)=x_1I+x_2J+y_1K+y_2L
\]
has determinant $\norm x^2-\norm y^2$, and the symmetric part
has eigenvalues $x_1\pm\norm y$. Thus the seed cone is
$\widehat C_I$. Under $M\mapsto AMB^T$, its image is
$\sym(BA^{-1}M)\succeq0$, so $F^T=BA^{-1}$ and
$\det F=1/(\det A)^2>0$. Transposition fixes the seed.
Conversely choose $A=\gamma I$, $B=\gamma F^T$, with
$\gamma^4\det F=1$, to obtain every $F$ of positive determinant.

The map $M\mapsto\sym(F^TM)$ is onto the symmetric matrices.
Writing $F^TM=H+tJ$,
$\det(F^TM)=\det H+t^2$. Hence positive definite symmetric
parts have positive determinant, and the whole PSD preimage
has nonnegative determinant. This proves full dimension and
freeness for $\det F>0$; a negative determinant reverses the
sign on the interior and is not free.
For a nonzero rank-one $ab^T$, the symmetric part of
$(F^Ta)b^T$ is PSD exactly when $F^Ta$ is a nonnegative
multiple of $b$. Since $F$ is invertible, that multiple is
positive for nonzero $a,b$.
With $a=(x,1)^T$, $b=(y,1)^T$, this gives the stated graph
and positive denominator. Its derivative is
$\det F/(F_{12}x+F_{22})^2>0$.
\end{proof}

\begin{proposition}[The BCM rotation family]\label{fd:bcm}
For $M=\bigl(\begin{smallmatrix}a&b\\c&d\end{smallmatrix}\bigr)$ and
$G_t=\bigl(\begin{smallmatrix}\cos t&\sin t\\-\sin t&\cos t\end{smallmatrix}\bigr)$,
\[
 \sym(G_tM)\succeq0
 \quad\Longleftrightarrow\quad
 \cos t(a+d)-\sin t(b-c)\ge\norm{(b+c,a-d)}.
\]
Thus BCM's (14a) is exactly the rotation subfamily.
If $F$ is not a scalar multiple of a rotation, its orbit cone
is no (14a) member. A positive-determinant orbit cone is
contained in no (14b) member.
\end{proposition}
\begin{proof}
The PSD trace criterion for a symmetric $2\times2$ matrix gives
the displayed identity: its traceless-part norm is unchanged
by the rotation. For $A=F^T$, the lineality of the cone
$\sym(AM)\succeq0$ is $\operatorname{span}(A^{-1}J)$.
Equality with a rotation member would make
$A^{-1}J$ proportional to $G_t^{-1}J$, forcing $A$ to be
a scalar multiple of $G_t$.
The interior of a positive-determinant orbit cone has
$ad-bc>0$. The (14b) inequality instead implies
$\norm{(b+c,a-d)}\ge\norm{(a+d,b-c)}$, that is $ad-bc\le0$.
These facts exclude containment. For example
$F^T=\diag(2,1/2)$ is outside the rotation family.
\end{proof}

\begin{proof}[Proof of \cref{fd:completion}]
\label{fd:completion-proof}
The normal indexed by $v$ is $Fvv^T$. Its null tangency matrix is
\[
 F^{-T}(Jv)(Jv)^T.
\]
The identity
$F^{-T}=JFJ^T/\det F$ gives its lower-right coordinate as
$(Fv)_1v_1/\det F$. Its sign therefore agrees with
$\ip{Fvv^T}{E}=v^TZ_Fv$.

The dual cone of $\widehat C_F$ is $F\mathbb S_+^2$.
The dual of its upward sum consists of $FY$ with
$Y\succeq0$ and $\ip Y{Z_F}\ge0$.
Such a $Y=LL^T$ decomposes into two rank-one kept terms.
For $u_\theta=(\cos\theta,\sin\theta)$, put
$f(\theta)=(Lu_\theta)^TZ_F(Lu_\theta)$.
Then $f(\theta)+f(\theta+\pi/2)=\ip Y{Z_F}\ge0$,
and the difference changes sign as $\theta$ increases by
$\pi/2$. At an intermediate angle both terms are equal
and nonnegative. The two corresponding rank-one terms
sum to $Y$. This works also for singular $L$.
The dual is therefore the cone generated by the kept
rank-one normals, proving by bipolarity the homogeneous
closed-sum identity.

Closure commutes with this positive slice for a specific reason.
If $M_n=A_n+\tau_nE\to M$ with $A_n\in\widehat C_F$,
$\tau_n\ge0$, and $h(M)>0$, then
\[
 0\le\det A_n=\det M_n-\tau_n h(M_n)
\]
bounds $\tau_n$. Passing to a subsequence gives a finite
lowering and $M-\tau E\in\widehat C_F$.
On $h=1$ it satisfies $0\le\tau\le q(s)$, proving
\cref{fd:lowering} and the slice identity. Determinant
nonnegativity also makes the completed slice a subset
of $\{q\ge0\}$, hence free when full dimensional.

Now $Z_F=\bigl(\begin{smallmatrix}a&b/2\\b/2&0\end{smallmatrix}\bigr)$.
If $b=0,a<0$, only vectors with first coordinate zero
are retained. Since $d<0$, their inequality is $h\le0$,
so the slice is empty. Otherwise $Z_F$ has positive
directions, dense in its nonnegative directions.
Each strictly kept normal has a tangency ray with $h>0$.
Scale its contact to $h=1$ and perturb its rank-one PSD
image to be positive definite; this proves that the
original cone and its completion meet the slice in
full dimension.

At such a contact $M_*=a_0b_0^T$, write
$F^TM_*=c\,b_0b_0^T$, $c>0$.
A dual normal $FY$ active there satisfies
$b_0^TYb_0=0$, forcing $Y$ onto the single rank-one
ray orthogonal to $b_0$. Thus the completed cone's
normal cone at the contact is that one ray.
The same is true in the slice after restricting the
normal, since the cone's interior meets the slice.
Any outside point violates a kept inequality and, by
density, a strictly kept one. Adding it makes the
associated contact interior to the convex enlargement:
there is then no nonzero supporting normal at that
contact. This contradicts freeness and proves maximality.
\end{proof}

\begin{remark}[The retained-cap support formula]\label{fd:scip-cap}
Let $G=\{\beta\in\R^m:\norm\beta=1,\ \beta_{\rm last}\le\gamma\}$,
where $m\ge2$ and $-1\le\gamma\le1$. For $y=(\widetilde y,y_{\rm last})$,
its support function is
\[
 \max_{\beta\in G}\beta^Ty=
 \begin{cases}
 \norm y,&y_{\rm last}\le\gamma\norm y,\\
 \sqrt{1-\gamma^2}\norm{\widetilde y}+\gamma y_{\rm last},
       &y_{\rm last}>\gamma\norm y .
 \end{cases}
\]
Indeed fixing $\beta_{\rm last}=t$ reduces the maximization to
$t y_{\rm last}+\sqrt{1-t^2}\norm{\widetilde y}$ for $t\le\gamma$.
The unconstrained optimum is $t=y_{\rm last}/\norm y$;
otherwise the endpoint $\gamma$ is optimal. The zero vector
and endpoint caps follow by continuity. This is the two-branch
function in bilinear Case~4 of \citet{ChmielaMunozSerrano2020ZIB}, with
$\gamma$ their last positive direction component. It explains
the identification with the constant-direction retained set.
\end{remark}

\subsection{The pencil and the exact starting examples}

\begin{proof}[Proof of \cref{fd:tangent-pencil}]
\label{fd:pencil-proof}
Let $M_*=M(t)=a_0b_0^T$, with
$a_0=(x_t,1)^T$ and $b_0=(y_t,1)^T$, and let $D=M_0(d)$.
At the contact no lowering is possible, so
$F^Ta_0=c b_0$, $c>0$.
For small positive and negative $s$, membership of $t+sd$
gives a lowering $0\le\tau(s)\le s^2\det D$.
Set $Y=\sym(F^TD)$, $Z=\sym(F^TE)$, and $v=Jb_0$.
The contact inequality is strictly kept, so $v^TZv>0$.
The null quadratic entry of the lowered PSD matrix is
\[
 s\,v^TYv-\tau(s)v^TZv\ge0.
\]
The $O(s^2)$ lowering bound for both signs of $s$
forces $v^TYv=0$, and then forces $\tau(s)=0$.
A PSD matrix with this zero diagonal entry has zero
off-diagonal entry in the basis $(b_0,v)$.
Hence $Y=\mu b_0b_0^T$ and
$\sym(F^T[D-(\mu/c)M_*])=0$.
Put $u=-\mu/c$. The matrix $D+uM_*$ has determinant
$\det D$, because $\det M_*=0$ and the mixed term
is $\nabla q(t)^Td=0$. It is invertible, so
$F^T(D+uM_*)=\theta'J$ and
$F^T=\theta J\adj(D+uM_*)$.
For $G_0=J\adj D$,
$G_0a_0=d_xb_0$; for $G_1=J\adj M_*$,
$G_1a_0=0$. Thus $\theta d_x>0$, and the chosen
$d_x>0$ gives $\theta>0$.
The determinant and $G_1M_*$ identities follow immediately.

If $\det D=0$, the same argument gives
$F^T(D+uM_*)=\theta'J$ with a singular left-hand side.
Thus $\theta'=0$, so $D+uM_*=0$.
Its lower-right entry gives $u=0$, and then $D=0$,
contrary to the segment being nontrivial.
\end{proof}

\begin{proof}[Proof of \cref{fd:counterexample}]
\label{fd:counterexample-proof}
Let $(\alpha,\beta,\gamma,\delta)$ be nonnegative
barycentric coordinates of $(\sbar,v_1,v_2,v_3)$.
For homogeneous barycentric coordinates, put
$Q=w(\alpha+\beta+\gamma+\delta)-xy$.
Direct expansion gives the exact identity
\begin{align*}
64Q={}&3[16(\beta-\gamma)-5\alpha+2\delta]^2
       +21\alpha^2+148\delta^2+1036\alpha\delta\\
     &+192\alpha\gamma+1280\beta\delta.
\end{align*}
On the simplex the coordinate sum is one.
Equality forces $\alpha=\delta=0$ and $\beta=\gamma=1/2$.
Thus $T_1\cap S=\{t\}$, and the invertible ray matrix
(determinant $249$) gives the unique corner minimizer
$(1/2,1/2,0)$ and value one.

For $d=v_1-v_2=(4,-12,36)$ the pencil is
\[
 G(u)=\begin{pmatrix}12&36\\-u&4-3u\end{pmatrix},
 \qquad \det G(u)=48.
\]
At the four vertices the PSD parameter intervals are
\[
 \begin{array}{c|c}
 \sbar &[12-8\sqrt2,\,12+8\sqrt2]\\
 v_1 &(-\infty,2]\\
 v_2 &[-2,\infty)\\
 v_3 &[-13/5-2\sqrt{30}/5,\,-13/5+2\sqrt{30}/5].
 \end{array}
\]
They follow by the two diagonal and determinant tests.
The first and last intervals are disjoint, excluding
all orbit slices containing $T_1$.

For completions set
$u_s=(1-\sqrt2/3,1)^T$,
$u_3=(-3/10+\sqrt{30}/60,1)^T$,
$r_s=12-8\sqrt2$, and
$r_3=-13/5+2\sqrt{30}/5<r_s$.
With $A_v(u)=\sym(G(u)M(v))$ and $Z(u)=\sym(G(u)E)$,
direct multiplication gives
\begin{align*}
 u_s^TA_{\sbar}(u)u_s&=(\sqrt2/2)(u-r_s),\\
 u_s^TZ(u)u_s&=-(1-\sqrt2/3)u+44/3-8\sqrt2>0
       &&(u\le r_s),\\
 u_3^TA_{v_3}(u)u_3&=-(\sqrt{30}/6)(u-r_3),\\
 u_3^TZ(u)u_3&=(3/10-\sqrt{30}/60)u+59/50-3\sqrt{30}/25>0
       &&(u\ge r_3).
\end{align*}
Thus no nonnegative lowering repairs the apex for $u<r_s$,
or the third vertex for $u>r_3$. Every real parameter is excluded.

To prove a strict supremum gap, suppose completions have
values tending to one. Normalize their parameters by
$\norm F_{\rm Frob}=1$ and choose feasible target simplices
$T_{z_n}$ with $z_n\to1$.
The vertex lowerings lie in $[0,q(v_{j,n})]$ and are bounded.
Pass to limits $F\ne0$ and $\tau_j$.
Every convex combination of the lowered limiting vertices
has nonnegative determinant: it is a limit of combinations
in the original homogeneous cones.
If $\det F>0$, the closed PSD conditions put $T_1$ in its
completion, contradicting the pencil certificate.

Otherwise $F$ has rank one. Write $F^T=ab^T$.
PSD of $\sym(a b^TM)$ forces $M^Tb$ into $\R_+a$,
so all lowered vertices lie in one affine plane $\Sigma$.
This is a genuine plane: the equation obtained by taking
a vector orthogonal to $a$ cannot vanish identically
on the entire $h=1$ slice for nonzero $a,b$.
The lowered midpoint of $v_1,v_2$ is
$M(t)-(\tau_1+\tau_2)E/2$; its determinant is
$-(\tau_1+\tau_2)/2$. Nonnegativity forces
$\tau_1=\tau_2=0$, so $t\in\Sigma$.
If $\Sigma$ is vertical, all original vertices also
lie in it, contradicting their affine independence.
Otherwise it is $w=\ell(x,y)$ for an affine function
$\ell$. Convex combinations of the lowered vertices
and nonnegative $E$ directions form its epigraph
over the projection of $T_1$. Their determinants
are nonnegative, so $\ell(x,y)\ge xy$ throughout
that projection, with equality at $(-3,0)$.
This point is interior: its projection is the strict
convex combination with weights
$(184,287,257,72)/800$ of the four projected vertices.
The derivative at this interior minimum forces
$\ell$ to be the tangent plane to $xy$ there.
But $\ell-xy=-(x+3)y$ changes sign in every
neighbourhood of that point. This contradiction
excludes rank-one limits and proves $\zB<1$.

Finally, $\nabla q(t)^T(\sbar-t)=3/2>0$, so
\cref{fd:attainment} gives unrestricted attainment.
For stability, positive costs make nearby bounded
objective sublevels uniformly compact. Limits of
minimizers give lower semicontinuity of the corner
value. A small radial extension of the original
minimizer has $q<0$, giving upper semicontinuity.
The value is therefore continuous here, and every
nearby minimizer approaches the unique original one.
The strict transversality value consequently persists.
If no neighbourhood preserved the strict family gap,
choose approaching data and family cuts with values
approaching the nearby corner values. The same
normalized-parameter and bounded-lowering limit
would produce one of the two contradictions just
proved at the original data. Thus the gap and
unrestricted attainment persist on an open neighbourhood.
\end{proof}

\begin{example}[A second rational orbit obstruction]
\label{fd:second-example}
Take costs $\mathbf1$, apex $(-1/2,1/2,1)$, and
\[
 v_1=(1/2,-4,-1/2),\qquad
 v_2=(-10,3,24),\qquad v_3=(9/2,1/2,15/2).
\]
The unique corner contact is $t=(-1,-3,3)$, with
ray vector $(6/7,1/7,0)$ and value one.
Here $\zA<1$. The statement asserts no certified
strict gap for the unrestricted completion family.
\end{example}
\begin{proof}
Let vertex zero denote the apex. On edge $ij$ use
$(1-\theta)v_i+\theta v_j$, $0\le\theta\le1$.
The exact edge polynomials and minima are
\[
 \begin{array}{c|c|c}
 ij&q(\theta)&\min q\\ \hline
 01&9\theta^2/2-17\theta/4+5/4&71/288\\
 02&95\theta^2/4+29\theta+5/4&5/4\\
 03&4\theta+5/4&5/4\\
 12&147\theta^2/2-21\theta+3/2&0\\
 13&-18\theta^2+87\theta/4+3/2&3/2\\
 23&145\theta^2/4-85\theta+54&21/4 .
 \end{array}
\]
The zero is unique on edge 12 at $\theta=1/7$.
The projected determinants of faces 012,013,023,123
are respectively $-161/4,45/2,-25/2,-301/4$.
Every triangular face therefore has independent
$x,y$ directions, so the restricted quadratic
Hessian is indefinite and cannot have an interior
minimum. The full tetrahedron also has an
indefinite Hessian. Thus all minima occur on
edges, proving nonnegativity and uniqueness.
The ray matrix is invertible, so the corner
minimizer is unique.

The tangent pencil gives the intervals
\[
 K_A(\sbar)=
 [133/12-7\sqrt{30}/6,\,133/12+7\sqrt{30}/6],
 \quad
 K_A(v_3)=
 [-1/14-3\sqrt{14}/7,\,-1/14+3\sqrt{14}/7].
\]
Their endpoints follow by the diagonal and
determinant tests for
$G(u)=\bigl(\begin{smallmatrix}3u+7&3u-49/2\\-u&21/2-u\end{smallmatrix}\bigr)$.
They are disjoint, since the lower endpoint
of the first exceeds $4$ whereas the upper
endpoint of the second is less than $2$.
Thus no invertible orbit parameter contains
$T_1$. If values approached one, normalized
parameters would have a nonzero limit.
A positive-determinant limit is excluded
by the pencil. A rank-one limit forces the
four vertices into a plane by the same
rank-one PSD identity used above, whereas
the ray determinant is $-6083/8\ne0$.
This proves $\zA<1$.
\end{proof}

\begin{proof}[Proof of \cref{fd:support-one}]
\label{fd:support-one-proof}
For the stated four vertices, homogeneous
barycentric coordinates give the identity
\begin{align*}
512Q={}&(64\alpha-79\gamma-32\delta)^2
       +31\gamma^2+512\delta^2+1280\alpha\delta
       +4032\gamma\delta\\
      &+128\beta(8\alpha+\gamma+2\delta).
\end{align*}
Equality forces $\alpha=\gamma=\delta=0$ and
$\beta=1$. Hence $T_1\cap S=\{v_1\}$.
The ray determinant is $67/2$, so the minimizer
is uniquely $e_1$.
At zero $\nabla q=(0,0,1)$, giving edge derivatives
$2,1/4,1/2$. Also
$P^T\nabla q=(-2,-7/4,-3/2)$.
The KKT multiplier $\sigma=1/2$ gives the stated
inactive multipliers.

For vertices $(\sbar,v_1,v_2,v_3)$ respectively,
take the positive definite matrices
\[
 Y_0=\begin{pmatrix}92&-75\\-75&64\end{pmatrix},\
 Y_1=\begin{pmatrix}726&89\\89&12\end{pmatrix},\
 Y_2=\begin{pmatrix}96&-64\\-64&45\end{pmatrix},\
 Y_3=\begin{pmatrix}20&16\\16&16\end{pmatrix}.
\]
Their determinants are $263,791,224,64$, and
direct multiplication gives
\[
 \sum_{j=0}^3 M(v_j)Y_j
 =\begin{pmatrix}334&-278\\201&-161\end{pmatrix}
 +\begin{pmatrix}0&0\\89&12\end{pmatrix}
 +\begin{pmatrix}-360&254\\-256&173\end{pmatrix}
 +\begin{pmatrix}26&24\\-34&-24\end{pmatrix}=0.
\]
If every $\sym(F^TM(v_j))$ were PSD, pairing
with the $Y_j$ and summing would give zero.
Positive definiteness forces every PSD matrix
to be zero. The four affine-independent
$M(v_j)$ span all $2\times2$ matrices, and
$\sym(F^TM)=0$ for all matrices forces $F=0$.
Thus even a nonzero singular $F$ is excluded.
If orbit cut values approached one, a
Frobenius-normalized parameter sequence would
converge to a nonzero matrix satisfying these
closed conditions. This contradiction proves
the strict supremum gap.
\end{proof}

\begin{proof}[Proof of \cref{fd:wedge}]
\label{fd:wedge-proof}
Use the affine shear
$X=x-a$, $Y=y-b$, $W=w-(bx+ay-ab)$.
Then $q=W-XY$ and the wedge is
$X\ge0,Y\le0,W\ge0$, whose interior is free.
If an outside point has $X<0$, mix it with
an interior wedge point $(\epsilon,-T,1)$,
where $\epsilon$ is small and $T$ large.
The resulting interior point of the convex
enlargement has $XY$ arbitrarily large
positive while $W$ stays bounded, hence $q<0$.
For $Y>0$ use $(T,-\epsilon,1)$.
The remaining outside points have $X\ge0,Y\le0$
and $W<0$. Mix a weight $\epsilon$ of such a
point with $(\epsilon^2,-\epsilon^2,\epsilon^2)$.
Its $W$ coordinate is negative of order
$\epsilon$, whereas $XY$ is of order
$\epsilon^2$, again giving $q<0$ in the
enlargement's interior. This proves maximality.

At any contact $q=0$, a completion lowering
must be zero. Two distinct contacts with the
same $x$ have $M_i=a b_i^T$ with the same
nonzero $a$ and independent $b_i$.
The rank-one membership identity would force
$F^Ta\in\R_+b_1\cap\R_+b_2=\{0\}$, impossible
for invertible $F$. The same applies with
$y$ fixed. The wedge contains such contacts,
so neither family contains it.
\end{proof}
